\begin{table}[htbp]
\centering
\caption{Spot Returns and Net Foreign Assets}
\label{tab:nfa_panel}
\begin{tabular}{lcccc}
\hline\hline
 & (1) & (2) & (3) & (4) \\
 & Main & Two-Way & Country & MP1 \\
 & (Date Cl.) & (Two-Way Cl.) & (Country Cl.) & (Date Cl.) \\
\hline
STMT & -0.5978 & -0.5978 & -0.5978*** & \\
 & (0.6176) & (0.4953) & (0.2160) & \\
MP1 & & & & 0.0216 \\
 & & & & (0.4489) \\
Surprise $\times$ NFA/GDP & -0.0048 & -0.0048 & -0.0048 & -0.0162* \\
 & (0.0067) & (0.0036) & (0.0051) & (0.0091) \\
NFA/GDP & 0.0004 & 0.0004** & 0.0004*** & 0.0003 \\
 & (0.0004) & (0.0002) & (0.0001) & (0.0004) \\
\hline
Country FE & Yes & Yes & Yes & Yes \\
Date Cluster & Yes & Yes & No & Yes \\
Country Cluster & No & Yes & Yes & No \\
Observations & 2103 & 2103 & 2103 & 2103 \\
$R^2$ & 0.0014 & 0.0014 & 0.0014 & 0.0059 \\
\hline\hline
\end{tabular}
\begin{tablenotes}
\small
\item \textit{Notes:} The dependent variable is the spot return $r_{i,t} = -\Delta\log(e_{i,t})$,
where $e$ is the foreign-per-USD exchange rate. Positive $r$ denotes foreign currency
appreciation. This convention matches Antol\'{i}n-D\'{i}az et al.\ (2023), so that
$\beta_2 > 0$ means creditor countries depreciate less. STMT is the statement surprise;
MP1 is the target rate surprise, both in basis points. NFA/GDP is from Lane \&
Milesi-Ferretti (EWN, 2024 update), lagged one year and demeaned so $\beta_1$
represents the effect at average NFA.
\textbf{Standard errors:} Column (1) clusters by FOMC date, the preferred
specification because STMT$_t$ is identical across currencies on each date,
inducing cross-sectional correlation \citep{Petersen2009}. Column (2) clusters
two-way (country $+$ date) as a conservative robustness check. Column (3)
clusters by country only; while reported for completeness, it is inappropriate
as it ignores the common-shock structure.
Sample: 8 currencies (AUD, CAD, CHF, EUR, GBP, JPY, MXN, NOK)
$\times$ 269 FOMC events, 1994--2024.
*** $p<0.01$, ** $p<0.05$, * $p<0.10$.
\end{tablenotes}
\end{table}
